\subsection{SEPARATION OF SOUSLIN SETS}

THEOREM 2. Let X be a metrizable space. If we are given a sequence \((\mathbf{X}_{n})\) of mutually disjoint Souslin subspaces of X, then there exists a sequence \((\mathbf{B}_{n})\) of mutually disjoint Borel sets in X such that \(X_{n} \subset B_{n}\) for each n.

The proof rests on two lemmas :

LEMMA 4. Let \((\mathbf{A}_n), (\mathbf{A}_m')\) be two sequences of subsets of a topological space \(\mathbf{X}\) . Suppose that for each pair \((\mathbf{A}_n, \mathbf{A}_m')\) there is a Borel set \(\mathbf{B}_{nm}\) of \(\mathbf{X}\) such that \(\mathbf{B}_{nm} \supset \mathbf{A}_n\) and \(\mathbf{B}_{nm} \cap \mathbf{A}_m' = \emptyset\) . Then there is a Borel set \(\mathbf{B}\) of \(\mathbf{X}\) which contains \(\bigcup_{n} \mathbf{A}_n\) and does not meet \(\bigcup_{m} \mathbf{A}_m'\) .

For the Borel set \(\mathbf{B} = \bigcup_{n}\left(\bigcap_{m}\mathrm{B}_{nm}\right)\) satisfies these conditions.

LEMMA 5. Let X be a Hausdorff space and let A, A' be two disjoint Souslin subspaces of X. Then there is a Borel set B of X such that B ⊃ A and B ∩ A' = ∅.

By Proposition 13 of no. 5 there exist two sieves \(\mathbf{C}\) , \(\mathbf{C}'\) and continuous mappings \(f\) of \(\mathbf{L}(\mathbf{C})\) onto \(\mathbf{A}\) , \(f'\) of \(\mathbf{L}(\mathbf{C}')\) onto \(\mathbf{A}'\) , constructed by the method described in no. 5. For each \(n \geqslant 0\) and each \(c \in \mathbf{C}_n\) , let \(q_n(c)\) denote the subspace of \(\mathbf{L}(\mathbf{C})\) formed of all sequences \((c_k)_{k \geqslant 0}\) such that \(c_n = c\) ; \(q_n(c)\) is a closed subspace of \(\mathbf{L}(\mathbf{C})\) . For each \(\gamma = (c_n) \in \mathbf{L}(\mathbf{C})\) the sequence of closed sets \(q_n(c_n)\) is decreasing and forms a filter base with \(\gamma\) as limit. Moreover, for each \(c \in \mathbf{C}_n\) , the sets \(q_{n+1}(d)\) , where \(d\) runs through the set \(\overline{p_n}(c)\) in \(\mathbf{C}_{n+1}\) , form a partition of \(q_n(c)\) . Make the analogous definitions of notation for the sieve \(c'\) .

We shall argue by contradiction and assume that every Borel set which contains A meets \(A'\) . In the first place, it follows from Lemma 4 and the definition of a sifting that there exist \(c_{0} \in C_{0}\) and \(c_{0}' \in C_{0}'\) such that every Borel set containing \(f(q_{0}(c_{0}))\) meets \(f'(q_{0}'(c_{0}'))\) . We can then define, by induction on n,

\[
\gamma = (c _ {n}) \in \mathrm{L} (\mathrm{C}) \quad \text { and } \quad \gamma^ {\prime} = (c _ {n} ^ {\prime}) \in \mathrm{L} (\mathrm{C} ^ {\prime})
\]

as follows: suppose that \(c_i\) and \(c_i'\) have already been defined for \(i < n\) , in such a way that for each index \(i < n\) every Borel set containing \(f(q_i(c_i))\) meets \(f'(q_i'(c_i'))\) ; applying Lemma 4 and the definition of a sifting to the sets \(f(q_{n-1}(c_{n-1}))\) and \(f'(q_{n-1}'(c_{n-1}'))\) , we see that there exist \(c_n \in \mathbf{C}_n\) and \(c_n' \in \mathbf{C}_n'\) such that \(p_{n-1}(c_n) = c_{n-1}\) and \(p_{n-1}(c_n') = c_{n-1}'\) .

and such that every Borel set containing \(f(q_n(c_n))\) meets \(f'(q_n'(c_n'))\) . Now the sequence \(f(q_n(c_n))\) converges to a point \(a = f(\gamma) \in \mathbf{A}\) , and the sequence \(f'(q_n'(c_n'))\) converges to a point \(a' = f'(\gamma') \in \mathbf{A}'\) . Since \(\mathbf{A} \cap \mathbf{A}' = \emptyset\) and \(\mathbf{X}\) is Hausdorff, there is a closed neighbourhood \(\mathbf{V}\) of \(a\) which does not contain \(a'\) , and thus for \(n\) sufficiently large, \(\mathbf{V}\) contains \(f(q_n(c_n))\) and does not meet \(f'(q_n'(c_n'))\) . This is a contradiction, since \(\mathbf{V}\) is a Borel set.

To prove Theorem 2, let \(\mathbf{Y}_n\) denote the union of the sets \(\mathbf{X}_i\) such that \(i \neq n\) ; then \(\mathbf{Y}_n\) is a Souslin subspace of \(\mathbf{X}\) (no. 2, Proposition 8). For each index \(n\) there exists a Borel set \(\mathbf{B}_n'\) which contains \(\mathbf{X}_n\) and does not meet \(\mathbf{Y}_n\) , by Lemma 5. Let \(\mathbf{B}_n\) be the intersection of \(\mathbf{B}_n'\) and \(\bigcap_{i < n} (\mathbf{X} - \mathbf{B}_i')\) . Then the \(\mathbf{B}_n\) are Borel sets, are mutually disjoint, and are such that \(\mathbf{B}_n \supset \mathbf{X}_n\) for each \(n\) .

COROLLARY. If a countable partition of a metrizable space is formed of Souslin sets, then these sets are Borel sets. In particular, every Souslin set in a metrizable space, whose complement is a Souslin set, is a Borel set.

